% Spectres infrarouges des deux produits possibles de l'oxydation du propan-1-ol (exercice 12, chapitre 8) ; schématiques
\begin{tikzpicture}
\begin{axis}[otIR,width=72mm,name=a,title={Spectre A}]
\addplot[otSpectre] {100 \bandeIR{2980}{25}{40} \bandeIR{2900}{25}{20} \bandeIR{2820}{12}{28} \bandeIR{2720}{12}{32}
  \bandeIR{1730}{16}{85} \bandeIR{1460}{15}{25} \bandeIR{1390}{15}{20} \bandeIR{1110}{20}{30} \bandeIR{850}{15}{20}};
\solution{
  \node[pcLegende,text=otGreenDark,anchor=north] at (axis cs:2770,62) {C--H ald.};
  \node[pcLegende,text=otGreenDark,anchor=west] at (axis cs:1700,17) {C=O};
  \node[pcLegende,text=otGreenDark,anchor=south] at (axis cs:3600,6) {propanal};}
\end{axis}
\begin{axis}[otIR,width=72mm,name=b,at={(a.east)},anchor=west,xshift=12mm,title={Spectre B},ylabel={}]
\addplot[otSpectre] {100 \bandeIR{2980}{300}{45} \bandeIR{1712}{20}{80} \bandeIR{1420}{20}{25} \bandeIR{1240}{25}{45}
  \bandeIR{940}{25}{30}};
\solution{
  \node[pcLegende,text=otGreenDark] at (axis cs:2980,38) {O--H acide};
  \node[pcLegende,text=otGreenDark,anchor=west] at (axis cs:1690,22) {C=O};
  \node[pcLegende,text=otGreenDark,anchor=south] at (axis cs:3600,6) {acide propanoïque};}
\end{axis}
\end{tikzpicture}
